\documentclass[11pt,a4paper]{article}
\usepackage[utf8]{inputenc}
\usepackage[vietnamese]{babel}
\usepackage{amsmath,amssymb}
\usepackage{geometry}
\geometry{a4paper, total={180mm,260mm}, left=15mm, right=15mm, top=18mm, bottom=18mm}
\usepackage{tikz}
\usetikzlibrary{shapes.geometric, arrows.meta, positioning, calc}
\usepackage{tcolorbox}
\tcbuselibrary{skins,breakable}
\usepackage{xcolor}
\usepackage{multicol}
\usepackage{array}
\usepackage{fontawesome}
\usepackage{fancyhdr}

\pagestyle{fancy}
\fancyhf{}
\renewcommand{\headrulewidth}{0.5pt}
\renewcommand{\footrulewidth}{0.5pt}
\fancyhead[L]{\small\bfseries\color{teal}\faBook\ SỔ TAY KIẾN THỨC TOÁN 3}
\fancyhead[R]{\small\bfseries\color{teal}Lớp Toán Thầy Lĩnh -- loptoanthaylinh.vn}
\fancyfoot[C]{\small\color{gray}-- Sổ tay kiến thức: Bài 73. Thu thập, phân loại, ghi chép số liệu. Bảng số liệu --}

\definecolor{primary}{HTML}{005F73}
\definecolor{secondary}{HTML}{0A9396}
\definecolor{accent}{HTML}{EE9B00}
\definecolor{danger}{HTML}{AE2012}
\definecolor{boxbg}{HTML}{F4F6F8}

\begin{document}

% --- BANNER TIÊU ĐỀ ---
\begin{tcolorbox}[
  enhanced,
  colback=white,
  colframe=primary,
  arc=3mm,
  boxrule=1.2pt,
  left=4mm, right=4mm, top=2.5mm, bottom=2.5mm,
  drop shadow=black!10!white
]
  \centering
  {\Large\bfseries\color{primary}\faBookmark\ BÀI 73: THU THẬP, PHÂN LOẠI, BẢNG SỐ LIỆU}\\[1.5mm]
  {\small\itshape\color{black!75}Kĩ năng kiểm đếm, phân loại đối tượng, cấu trúc bảng số liệu (hàng, cột) và phân tích so sánh thông tin trên 1 trang A4}
\end{tcolorbox}

\vspace{-2mm}


\noindent
\begin{minipage}[t]{0.49\textwidth}\vspace{0pt}
\begin{tcolorbox}[
  enhanced,
  colback=blue!4!white,
  colframe=blue!75!black,
  arc=2mm, boxrule=0.8pt,
  left=2mm, right=2mm, top=1.8mm, bottom=1.8mm,
  title={\footnotesize\bfseries\color{white}\faKey\ 1. Thu thập, phân loại và ghi chép}
]
{\footnotesize
\begin{itemize}\setlength{\itemsep}{1.5pt}\setlength{\leftskip}{-2.5mm}
  \item \textbf{Thu thập số liệu:} Quan sát, đo đạc hoặc phỏng vấn để thu thập thông tin về một đối tượng.
  \item \textbf{Phân loại:} Nhóm các đối tượng có cùng đặc điểm lại với nhau (màu sắc, hình dạng, sở thích...).
  \item \textbf{Ghi chép số liệu:} Đếm số lượng của từng nhóm và ghi lại kết quả (dùng gạch đếm hoặc con số).
\end{itemize}
}
\end{tcolorbox}
\end{minipage}
\hfill
\begin{minipage}[t]{0.49\textwidth}\vspace{0pt}
\begin{tcolorbox}[
  enhanced,
  colback=teal!4!white,
  colframe=teal!75!black,
  arc=2mm, boxrule=0.8pt,
  left=2mm, right=2mm, top=1.8mm, bottom=1.8mm,
  title={\footnotesize\bfseries\color{white}\faCheckSquare\ 2. Cấu trúc và đọc Bảng số liệu}
]
{\footnotesize
\begin{itemize}\setlength{\itemsep}{1.5pt}\setlength{\leftskip}{-2.5mm}
  \item Bảng số liệu thường gồm các \textbf{hàng ngang} và các \textbf{cột dọc}.
  \item Hàng/cột đầu tiên ghi tên các loại đối tượng.
  \item Hàng/cột tiếp theo ghi số lượng tương ứng của từng đối tượng.
  \item \textbf{Kĩ năng phân tích bảng số liệu:}\\
  -- Nhìn ô đối tượng để đọc ngay số lượng.\\
  -- So sánh đối tượng nào nhiều nhất, ít nhất.\\
  -- Tính tổng số lượng hoặc phần chênh lệch.
\end{itemize}
}
\end{tcolorbox}
\end{minipage}

\vspace{-2.5mm}
% --- KHỐI 3: SƠ ĐỒ TƯ DUY TIKZ ---
\noindent
\begin{tcolorbox}[
  enhanced,
  colback=orange!4!white,
  colframe=orange!80!black,
  arc=2mm, boxrule=0.8pt,
  left=2mm, right=2mm, top=1.5mm, bottom=1.5mm,
  title={\footnotesize\bfseries\color{white}\faSitemap\ SƠ ĐỒ TƯ DUY: TIẾN TRÌNH 4 BƯỚC XỬ LÝ SỐ LIỆU THỐNG KÊ}
]
\centering
\begin{tikzpicture}[
  box/.style={draw, rounded corners=2pt, align=center, font=\scriptsize\bfseries, inner sep=2.5pt, minimum height=0.75cm},
  arr/.style={-{Stealth[length=2mm]}, thick}
]
  \node[box, fill=purple!15, draw=purple!80!black, text width=2.6cm] (s1) at (0, 0) {1. THU THẬP\\Quan sát, điều tra\\Thu nhận dữ liệu};
  \node[box, fill=teal!15, draw=teal!80!black, text width=2.6cm] (s2) at (3.5, 0) {2. PHÂN LOẠI\\Chia thành từng nhóm\\có đặc điểm chung};
  \node[box, fill=blue!12, draw=blue!80!black, text width=2.6cm] (s3) at (7.0, 0) {3. LẬP BẢNG\\Ghi chép số liệu\\vào các hàng, cột};
  \node[box, fill=orange!15, draw=orange!80!black, text width=2.6cm] (s4) at (10.5, 0) {4. PHÂN TÍCH\\So sánh nhiều/ít nhất\\Tính tổng và hiệu};

  \draw[arr, purple!80!black] (s1) -- (s2);
  \draw[arr, teal!80!black] (s2) -- (s3);
  \draw[arr, blue!80!black] (s3) -- (s4);
\end{tikzpicture}
\end{tcolorbox}

\vspace{-2.5mm}
% --- KHỐI 4: VÍ DỤ & MẸO TRÁNH BẪY ---
\noindent
\begin{minipage}[t]{0.49\textwidth}\vspace{0pt}
\begin{tcolorbox}[
  enhanced,
  colback=green!4!white,
  colframe=green!60!black,
  arc=2mm, boxrule=0.8pt,
  left=2mm, right=2mm, top=1.8mm, bottom=1.8mm,
  title={\footnotesize\bfseries\color{white}\faCheckCircle\ Ví dụ thực hành chuẩn mực}
]
{\footnotesize
\textbf{Bảng số cây hoa trồng trong vườn trường:}\\
\begin{center}
\renewcommand{\arraystretch}{1.1}
\begin{tabular}{|l|c|c|c|}
  \hline
  \textbf{Loài hoa} & Hoa hồng & Hoa cúc & Hoa đồng tiền \\
  \hline
  \textbf{Số lượng} & $25\text{ cây}$ & $40\text{ cây}$ & $15\text{ cây}$ \\
  \hline
\end{tabular}
\end{center}
-- Hoa cúc trồng \textbf{nhiều nhất} ($40$ cây).\\
-- Hoa đồng tiền trồng \textbf{ít nhất} ($15$ cây).\\
-- Tổng số cây hoa cả vườn là: $25 + 40 + 15 = \mathbf{80\text{ cây}}$.\\
-- Hoa cúc nhiều hơn hoa hồng: $40 - 25 = \mathbf{15\text{ cây}}$.
}
\end{tcolorbox}
\end{minipage}
\hfill
\begin{minipage}[t]{0.49\textwidth}\vspace{0pt}
\begin{tcolorbox}[
  enhanced,
  colback=red!4!white,
  colframe=red!70!black,
  arc=2mm, boxrule=0.8pt,
  left=2mm, right=2mm, top=1.8mm, bottom=1.8mm,
  title={\footnotesize\bfseries\color{white}\faExclamationTriangle\ Mẹo tránh bẫy trắc nghiệm}
]
{\footnotesize
\begin{itemize}\setlength{\itemsep}{1.5pt}\setlength{\leftskip}{-2.5mm}
  \item \textbf{Bẫy đọc nhầm dòng hoặc nhầm cột:}\\
  Khi tra cứu số liệu, hãy dùng ngón tay hoặc thước dóng thẳng hàng giữa tên đối tượng và con số tương ứng để không bị nhầm sang đối tượng kế bên!
  \item Đọc kĩ câu hỏi hỏi "nhiều nhất", "ít nhất" hay "tất cả bao nhiêu".
\end{itemize}
}
\end{tcolorbox}
\end{minipage}


\end{document}
